% Solution_Template.tex
\documentclass[11pt,a4paper]{article}
\usepackage[utf8]{inputenc}
\usepackage[T1]{fontenc}
\usepackage[ngerman, english]{babel} % remove ngerman if you prefer English only
\usepackage{geometry}
\usepackage{fancyhdr}
\usepackage{hyperref}
\usepackage{tikz}

\usetikzlibrary{trees}
\geometry{margin=25mm}
\pagestyle{fancy}
\fancyhf{}
\renewcommand{\headrulewidth}{0.4pt}
\lhead{\textbf{Name:} \studentname}
\rhead{\textbf{Matr.-Nr.:} \studentnumber}
\cfoot{\thepage}

% User info - edit these
\newcommand{\studentname}{Matthias Janßen}
\newcommand{\studentnumber}{1871808}
\newcommand{\course}{Übersetzung und Analyse von Programmen}
\newcommand{\institute}{Universität Trier}

\title{\course}
\author{\studentname\\Matrikelnummer: \studentnumber\\\institute}
\date{\today}

\begin{document}

\maketitle

\begin{center}
    \textbf{Abgabe:} Assignment / Sheet \\
    \vspace{4mm}
    \textbf{Betreuer:} Prof. Dr. Stephan Diehl\\Jan Niclas Ruppenthal, M.Sc.
\end{center}

\bigskip

\section*{Aufgaben}
% Start writing your solutions here.

\section{Aufgabe 1}
a) \\
    Terminale: \textit{T} = {\textit{a, e, g, l, r}}\\
    Nichtterminale: {S, A}    
\\   
    Produktionsregeln: \\

\begin{center}    
    S $\rightarrow \epsilon$ \\
    S $\rightarrow$ A \\

    A $\rightarrow$ aAa \\
    A $\rightarrow$ eAe \\
    A $\rightarrow$ gAg \\
    A $\rightarrow$ lAl \\
    A $\rightarrow$ rAr \\

    A $\rightarrow$ aa \\
    A $\rightarrow$ ee \\
    A $\rightarrow$ gg \\
    A $\rightarrow$ ll \\
    A $\rightarrow$ rr \\

    A $\rightarrow$ a \\
    A $\rightarrow$ e \\
    A $\rightarrow$ g \\
    A $\rightarrow$ l \\
    A $\rightarrow$ r \\

\end{center}

b)
Ableitungsbaum für das Wort "Lagerregal":
\begin{center}
    \begin{figure}
        \tikzset{
        % Set the overall layout of the tree
        level 1/.style = {level distance=3.5cm, sibling distance=3.5cm},
        level 2/.style = {level distance=3.5cm, sibling distance=2cm},
        % Define styles for bags and leafs
        bag/.style = {text width=4cm, text centered,  inner sep=1pt},
        end/.style = {circle, minimum width=3pt,fill, inner sep=0pt} 
        }
        \begin{tikzpicture}[grow=down]
            \tikzset{frontier/.style={distance from root=150pt}}
            
            \node {S}
                child {
                    node[bag] (l) {l}
                    edge from parent 
                    node[above] {}
                    node[below] {}
                }
                child {
                    node[bag] (A) {A}
                    child {
                        node[bag] (a) {a}
                        edge from parent
                        node[above] {}
                        node[below] {}
                    }
                    child {
                        node[bag] (a) {A}
                        child {
                            node[bag] (g) {g}
                            edge from parent
                            node[above] {}
                            node[below] {}
                        }
                        child {
                            node[bag] (g) {A}
                            child {
                                node[bag] (e) {e}
                                edge from parent
                                node[above] {}
                                node[below] {}
                            }
                            child {
                                node[bag] (e) {A}
                                child {
                                    node[bag] (r) {r}
                                    edge from parent
                                    node[above] {}
                                    node[below] {}
                                }
                                child {
                                    node[bag] (r) {r}
                                    edge from parent
                                    node[above] {}
                                    node[below] {}
                                }
                                edge from parent
                                node[above] {}
                                node[below] {}
                            }
                            child {
                                node[bag] (e) {e}
                                edge from parent
                                node[above] {}
                                node[below] {}
                            }
                            edge from parent
                            node[above] {}
                            node[below] {}
                        }
                        child {
                            node[bag] (g) {g}
                            edge from parent
                            node[above] {}
                            node[below] {}
                        }
                        edge from parent
                        node[above] {}
                        node[below] {}
                    }
                    child {
                        node[bag] (a) {a}
                        edge from parent
                        node[above] {}
                        node[below] {}
                    }
                    edge from parent 
                    node[above] {}
                    node[below] {}
                }
                    child {
                    node[bag] (l) {l}      
                    edge from parent 
                    node[above] {}
                    node[below] {}
                }
                ;
        \end{tikzpicture}
    \end{figure}
\end{center}    
\section{Aufgabe 2}
Nach einem gelesenen Zeichen gibt es einen Konflikt zwischen den Regeln 1 \& 2 da beide das Zeichen 
"if" enthalten. Nach Zeichen 2 ist der Konflikt weiterhin nicht gelöst da beide Regeln "(" zeigen.
Nach dem dritten Zeichen ist eindeutig, welche Regel angewendet wird, da Regel 1 "id" enthält während
in Regel 2 "not" ist. Die Grammatik ist daher LL(3).

\section{Aufgabe 3}

\section{Aufgabe 4}
a)
\begin{center}
    S $\rightarrow$ AB \\
    AB $\rightarrow$ BB (Regel: A $\rightarrow$ B) \\
    BB $\rightarrow$ CB (Regel: B $\rightarrow$ C) \\
    CB $\rightarrow$ AabB (Regel: C $\rightarrow$ Aab)
\end{center} 

Durch die richtige Reihenfolge der Regeln kann erneut das Nichtterminal A an der Linken Stelle erzeugt werden.
Die Grammatik ist daher nicht Rekursionsfrei.

b)
\begin{center}
    1.) S $\rightarrow$ AB \\
    AB $\rightarrow$ aBCB (Regel: A $\rightarrow$ aBC) \\
    aBCB $\rightarrow$ abCB (Regel: B $\rightarrow$ b) \\
    abCB $\rightarrow$ abB (Regel: C $\rightarrow$ $\epsilon$) \\
    abB $\rightarrow$ abC (Regel: B $\rightarrow$ C) \\
    abC $\rightarrow$ ab (Regel: C $\rightarrow$ $\epsilon$) \\
    \bigskip
    2.) S $\rightarrow$ AB \\
    AB $\rightarrow$ BB (Regel: AB $\rightarrow$ B)\\
    BB $\rightarrow$ aAB (Regel: B $\rightarrow$ aA)\\
    aAB $\rightarrow$ aBB (Regel: A $\rightarrow$ B)\\
    aBB $\rightarrow$ abB (Regel: B $\rightarrow$ b) \\
    abB $\rightarrow$ abC (Regel: B $\rightarrow$ C) \\
    abC $\rightarrow$ ab (Regel: C $\rightarrow$$ \epsilon$)
\end{center}

Es ist möglich auf verschiedene Arten das Wort "ab" zu erzeugen, damit ist die Grammatik mehrdeutig.


\end{document}